% Manuscript format standard/template
% ===================================
% Purpose:
%   Reusable LaTeX manuscript-format standard for journal-style .tex files.
%   Copy this file into another repository and replace the manuscript content.
%
% Reference style source:
%   /Users/yaolinjia/Desktop/Rice/paper/manuscript.tex
%
% Scope:
%   Defines title, author block, affiliations, abstract, keywords, headings,
%   body text, line spacing, fonts, font sizes, captions, figures, tables,
%   page numbers, line numbers, margins, hyperlinks, references, and common
%   manuscript helper commands.
%
% Compilation:
%   pdflatex manuscript_format_standard.tex
%   bibtex manuscript_format_standard
%   pdflatex manuscript_format_standard.tex
%   pdflatex manuscript_format_standard.tex
%
% Recommended use in a real paper:
%   1. Copy the preamble from this file into manuscript.tex.
%   2. Replace title, author block, abstract, keywords, sections, figures,
%      tables, and bibliography database name.
%   3. Keep the formatting commands unchanged unless required by the journal.

\documentclass[12pt]{article}

% -----------------------------------------------------------------------------
% Page size, margins, font encoding, and base fonts
% -----------------------------------------------------------------------------
\usepackage[margin=1in]{geometry}       % 1-inch page margins on all sides
\usepackage[T1]{fontenc}
\usepackage[utf8]{inputenc}
\usepackage{textcomp}
\usepackage{newtxtext,newtxmath}        % Times-like text and math fonts

% -----------------------------------------------------------------------------
% Core packages
% -----------------------------------------------------------------------------
\usepackage{amsmath}
\usepackage{graphicx}
\usepackage{booktabs}
\usepackage{array}
\usepackage{longtable}
\usepackage{calc}
\usepackage{float}
\usepackage{caption}
\usepackage{setspace}
\usepackage{lineno}
\usepackage{xcolor}
\usepackage[round,authoryear]{natbib}
\usepackage{hyperref}
\usepackage{bookmark}
\IfFileExists{xurl.sty}{\usepackage{xurl}}{}
\urlstyle{same}

% -----------------------------------------------------------------------------
% Page numbers, line spacing, paragraph spacing, and line breaking
% -----------------------------------------------------------------------------
\pagestyle{plain}                       % Page number centered in footer
\doublespacing                          % Double-spaced manuscript body
\setlength{\parindent}{0pt}             % No paragraph indentation
\setlength{\parskip}{0pt}               % No extra paragraph gap
\setlength{\emergencystretch}{3em}      % Reduce overfull lines

% -----------------------------------------------------------------------------
% Section heading format
% -----------------------------------------------------------------------------
% The reference manuscript uses standard article headings. The following sizes
% make the hierarchy explicit while preserving a conservative journal style.
\usepackage{titlesec}
\titleformat{\section}
  {\normalfont\large\bfseries}
  {\thesection}{0.75em}{}
\titleformat{\subsection}
  {\normalfont\normalsize\bfseries}
  {\thesubsection}{0.75em}{}
\titleformat{\subsubsection}
  {\normalfont\normalsize\itshape}
  {\thesubsubsection}{0.75em}{}
\titlespacing*{\section}{0pt}{1.2em}{0.6em}
\titlespacing*{\subsection}{0pt}{1.0em}{0.45em}
\titlespacing*{\subsubsection}{0pt}{0.8em}{0.35em}

% -----------------------------------------------------------------------------
% Caption and table font standards
% -----------------------------------------------------------------------------
\newcommand{\tablecaptionfontsize}{\fontsize{10pt}{12pt}\selectfont}
\newcommand{\tablefontsize}{\fontsize{9pt}{11pt}\selectfont}
\DeclareCaptionFont{captionten}{\fontsize{10pt}{12pt}\selectfont}
\captionsetup{
  font=captionten,
  labelfont=bf,
  textfont=captionten,
  justification=raggedright,
  labelsep=period,
  singlelinecheck=false
}
\newcolumntype{L}[1]{>{\tablefontsize\raggedright\arraybackslash}m{#1}}

% -----------------------------------------------------------------------------
% Hyperlink and bibliography helper commands
% -----------------------------------------------------------------------------
\providecommand{\tightlist}{\setlength{\itemsep}{0pt}\setlength{\parskip}{0pt}}
\providecommand{\DOIprefix}{doi: }
\providecommand{\URLprefix}{URL: }
\providecommand{\ArXivprefix}{arXiv: }
\providecommand{\Pubmedprefix}{PMID: }
\providecommand{\doi}[1]{\href{https://doi.org/\detokenize{#1}}{\nolinkurl{#1}}}
\hypersetup{
  pdftitle={Manuscript title},
  hidelinks,
  pdfcreator={LaTeX}
}

% Figure/table cross-reference helpers
\newcommand{\figref}[1]{\hyperref[#1]{Figure~\ref*{#1}}}
\newcommand{\figpanelref}[2]{\hyperref[#1]{Figure~\ref*{#1}#2}}
\newcommand{\tabref}[2]{\hyperref[#1]{Table~#2}}

% -----------------------------------------------------------------------------
% Manuscript title
% -----------------------------------------------------------------------------
% Title format:
%   - Centered
%   - Large font
%   - Slightly compressed line spacing for long titles
\newcommand{\manuscripttitle}{Manuscript title: replace with the full title}

\begin{document}

% -----------------------------------------------------------------------------
% Line numbers
% -----------------------------------------------------------------------------
% Enable line numbers for review manuscripts. Remove or comment this line for
% final production files if the journal requests no line numbering.
\linenumbers

% -----------------------------------------------------------------------------
% Title block
% -----------------------------------------------------------------------------
\begin{center}
{\LARGE\setstretch{0.9}\selectfont \manuscripttitle\par}
\end{center}

\vspace{0.8em}

% -----------------------------------------------------------------------------
% Author and affiliation block
% -----------------------------------------------------------------------------
% Standard format:
%   - Author names on one or more lines
%   - Affiliations listed below with superscript letters
%   - Corresponding author marked with *
\begingroup
\normalsize
\setstretch{1.05}
\noindent
\mbox{First~Author}\textsuperscript{a},
\mbox{Second~Author}\textsuperscript{b},
\mbox{Corresponding~Author}\textsuperscript{a,*}
\par
\vspace{0.7em}

\noindent
\textsuperscript{a}Affiliation A, City, Postal Code, Country
\par
\noindent
\textsuperscript{b}Affiliation B, City, Postal Code, Country
\par
\vspace{0.55em}

\noindent
\textsuperscript{*}Corresponding author:
Corresponding Author,
E-mail: \href{mailto:corresponding.author@example.com}
{corresponding.author@example.com}
\par
\endgroup

\vspace{1em}

% -----------------------------------------------------------------------------
% Abstract and keywords
% -----------------------------------------------------------------------------
% Abstract format:
%   - Unnumbered section heading
%   - Same body font as manuscript
%   - Double spacing inherited from document unless changed locally
\section*{Abstract}

Replace this paragraph with the manuscript abstract. The abstract should be a
single coherent block unless the target journal requires structured headings.
It should state the background, gap, objective, methods, main results, and
conclusion without citations unless explicitly permitted by the journal.

\noindent\textbf{Keywords:} keyword one; keyword two; keyword three; keyword four; keyword five

% -----------------------------------------------------------------------------
% Main manuscript body
% -----------------------------------------------------------------------------
% Heading hierarchy:
%   \section{}       = first-level heading
%   \subsection{}    = second-level heading
%   \subsubsection{} = third-level heading
% Avoid deeper heading levels in most journal manuscripts.

\section{Introduction}\label{introduction}

Replace with manuscript text. Body text is 12 pt Times-like font, double
spaced, with 1-inch margins, no paragraph indent, and line numbers enabled.
Use author-year citations with natbib, for example \citep{example2024}.

\section{Materials and methods}\label{materials-and-methods}

\subsection{Study design}\label{study-design}

Replace with method text.

\subsubsection{Detailed method component}\label{detailed-method-component}

Replace with detailed method text.

% -----------------------------------------------------------------------------
% Figure format
% -----------------------------------------------------------------------------
% Standard figure rules:
%   - Use figure environment
%   - Place with [htbp] unless exact placement is required
%   - Center figure
%   - Use width=\linewidth for full-width figures
%   - Caption below figure
%   - Label after caption
\begin{figure}[htbp]
\centering
\IfFileExists{example_figure.png}
  {\includegraphics[width=\linewidth]{example_figure.png}}
  {\fbox{\parbox[c][0.25\textheight][c]{0.9\linewidth}{\centering Example figure placeholder}}}
\caption{Replace with a concise figure caption. The caption should explain the
main message, define panels if present, and define abbreviations used in the
figure.}
\label{fig:example}
\end{figure}

% -----------------------------------------------------------------------------
% Table format
% -----------------------------------------------------------------------------
% Standard table rules:
%   - Use booktabs rules: \toprule, \midrule, \bottomrule
%   - Avoid vertical lines
%   - Use \tablefontsize for dense tables
%   - Use L{} columns for wrapped text
%   - Caption above table
\begin{table}[H]
\centering
\caption{Replace with a concise table caption.}
\label{tab:example}
\tablefontsize
\renewcommand{\arraystretch}{1}
\begin{tabular}{@{}L{0.28\linewidth} L{0.32\linewidth} L{0.30\linewidth}@{}}
\toprule
Column 1 & Column 2 & Column 3 \\
\midrule
Example row & Wrapped text is supported in L columns & Example value \\
\bottomrule
\end{tabular}
\end{table}

\section{Results}\label{results}

Replace with results text. Present results in a logical order that matches the
study objectives. Avoid mixing interpretation-heavy discussion into this
section unless the target journal permits combined Results and Discussion.

\section{Discussion}\label{discussion}

Replace with discussion text. The discussion should interpret the results,
connect them to the stated gap, compare with prior work, define implications,
and state limitations.

\section{Conclusion}\label{conclusion}

Replace with the conclusion. Keep it concise and avoid introducing new
results.

% -----------------------------------------------------------------------------
% Optional end sections
% -----------------------------------------------------------------------------
\section*{Acknowledgements}

Replace with acknowledgements or remove this section.

\section*{Declaration of competing interest}

The authors declare that they have no known competing financial interests or
personal relationships that could have appeared to influence the work reported
in this paper.

\section*{Data and code availability}

Data and code availability statement.

% -----------------------------------------------------------------------------
% References
% -----------------------------------------------------------------------------
% Bibliography format:
%   - Author-year citation style through natbib
%   - Elsevier Harvard style by default, matching the reference manuscript
%   - Change references to the project-specific .bib filename if needed
\bibliographystyle{elsarticle-harv}
\bibliography{references}

\end{document}
